\section{Retained mass and a lower bound for fine sausages}
\label{retain:section}

We now turn the uniform mesh tail into a positive measure with the
required small-ball bound.  A subsequential weak limit of the retained densities will give the
lower Hausdorff bound. A separate prelimit estimate at the end of the
section supplies the sausage mass used by the ribbon construction.

Use the parameters and density $W_e$ of Section~\ref{mesh:section}.
The gauge $h$ is the one in \eqref{eq:gauge}. For each fixed $L\ge1$,
$h(Lr)/h(r)\to L^d$ as $r\downarrow0$, so these values differ by
bounded factors at small scales. This permits conformal changes of
coordinates on compact interior sets.

\subsection{A positive first moment before a finite time}

Fix a closed disk $K\Subset\mathbb H$ with a fixed free collar.  Define
\[
 \nu_e^T(dz)=\mathbf1_K(z)W_e(z)
       \mathbf1_{\{\tau_e(z)\le T\}}\,dA(z).
\]
Let $\delta_K=\operatorname{dist}(K,\partial\mathbb H)>0$.
Use the enlargement of $K$ by $\delta_K/2$ in the scaled tail
\eqref{mesh:normalized-tail}, so its permitted final distances include
$0<e<\delta_K/2$.  Choose $0<e_0\le\delta_K/(8\mathrm e)$, which
lies in this range and in the range of Lemma~\ref{ll:approach} for
every $z\in K$.
We first show that a deterministic finite $T$ can be chosen so that
\begin{equation}
 \inf_{0<e<e_0}\mathbb E\nu_e^T(K)\ge m_0>0,
 \qquad
 \sup_{0<e<e_0}\mathbb E\bigl[\nu_e^T(K)^2\bigr]\le m_2<\infty.
 \label{retain:first-second}
\end{equation}
The second bound follows by integrating the uniform upper tail
\eqref{mesh:normalized-tail}, after scaling to $K$.  We justify the first bound,
including the finite physical time needed for a tilted approach.

For an integer $m\ge1$, let $\rho_m(z)$ be the first time the live
point has conformal radius $R_0(z)e^{-m}$, and write $\mathbb P_z^{*,m}$ for
the finite centered law through $\rho_m(z)$ from
Lemma~\ref{ll:finite-tilt}.  If $R_0(z)e^{-m}\le e$, the distance stop
$\tau_e(z)$ occurs by $\rho_m(z)$ under this law: Koebe's inequality
gives distance to the boundary at most the conformal radius, and the
initial boundary is farther than $e$.  The finite stopped density
identity therefore gives
\[
 \mathbb E\!\left[W_e(z)\mathbf1_{\{\tau_e(z)\le T\}}\right]
 =M_0(z)\mathbb P_z^{*,m}\{\tau_e(z)\le T\}
 \ge M_0(z)\mathbb P_z^{*,m}\{\rho_m(z)\le T\}.
\]
By consistency, the probabilities on the right decrease with $m$.
Define their limit by
\[
 q_T(z)=\lim_{m\to\infty}
       \mathbb P_z^{*,m}\{\rho_m(z)\le T\}.
\]
By continuity from above on the consistent centered path, this is
the probability that $T^*(z)=\lim_m\rho_m(z)$ is at most $T$.
Lemma~\ref{ll:finite-clock} therefore gives $q_T(z)\uparrow1$ as
$T\to\infty$.  Here the initial uniformization is the identity, so
the lemma's Loewner time is the original capacity time.
Each probability under a finite stopped law is a
measurable kernel in $z$ by the localized construction, so the same
is true of $q_T$.  Taking the decreasing
limit above and integrating over $K$ yields
\[
 \mathbb E\nu_e^T(K)\ge\int_K M_0(z)q_T(z)\,dA(z).
\]
The right side increases to $\int_KM_0(z)dA(z)>0$ as $T\to\infty$.
This proves \eqref{retain:first-second} for a sufficiently large
deterministic $T$, using density identities only at finite stops.

Brownian scaling places the same experiment in every prescribed
initial capacity interval of length $\Delta>0$.  Scaling space by
$c=\sqrt{\Delta/T}$ scales capacity time by $c^2$, the permitted mesh
cutoff by $c$, and the integrated density by $c^d$.  Thus the first
moment lower bound becomes $c^dm_0$ and the second moment upper bound
becomes $c^{2d}m_2$, leaving the ratio $m_0^2/m_2$ unchanged.

\subsection{Removing squares with excessive mass}

Take the deterministic sequence $e_N=2^{-N}$ and write
$\nu_N=\nu_{e_N}^T$.  Dyadic squares are taken half open and form a
partition at each level.  Let $k_0$ be a starting level to be fixed below.  For a
square $Q$ of side $l=2^{-k}$, $k_0\le k\le N$, declare it excessive
when
\[
 \nu_N(Q)>H_k,
 \qquad H_k=D l^d(\log k)^{p/2},
\]
where $D$ will be chosen large.  We claim that
\begin{equation}
 \mathbb E\left[\nu_N(Q)\mathbf1_{\{\nu_N(Q)>H_k\}}\right]
 \le C l^{p+d}k^{-\beta},
 \qquad \beta=cD^{2/p},
 \label{retain:bad-square}
\end{equation}
after reducing the positive constant $c$ if necessary.  Constants
are uniform in $N\ge k$ and in the squares meeting $K$.

Let $z_Q$ be the center of $Q$, let $\sigma_Q$ be the first hit of
$\partial B(z_Q,10l)$, and set
\[
 E_Q=\{\sigma_Q<\infty,\ z_Q\in D_{\sigma_Q}\}.
\]
Require $k_0$ to be large enough that $2^{-k_0}<e_0$ and
$10l\le\operatorname{dist}(z_Q,\partial\mathbb H)/(8\mathrm e)$
for every square under consideration.  In particular every
$\overline B(z_Q,10l)$ under consideration is initially interior.
We first check that only $E_Q$ can contribute to $\nu_N(Q)$.
Since $Q\subset B(z_Q,l)$ and $e_N\le l$, such a contribution requires
the trace to enter $B(z_Q,2l)$ after first hitting the larger circle.
If $z_Q$ is swallowed before or at that first hit, the open larger
disk is still disjoint from the trace at the swallowing time.  The
connected-component observation after Lemma~\ref{geom:first-hit-access}
then shows that its whole interior is swallowed, so no later trace
can enter it.
Thus $\nu_N(Q)>0$ implies $E_Q$.

On $E_Q$, the hit is a first live approach to $z_Q$ at distance $10l$.
Lemma~\ref{ll:approach} therefore gives
$\mathbb P(E_Q)\le Cl^p$, uniformly because $M_0(z_Q)$ is bounded on
a fixed compact neighborhood of $K$.  Also the open disk
$B(z_Q,10l)$ is contained in $D_{\sigma_Q}$: it is trace-free,
connected, and contains the live center.  In particular
$A_Q=B(z_Q,2l)$ contains $Q$ and its closed enlargement
$\overline B(z_Q,6l)$ remains free.

Conditionally on $E_Q$ and the stopped past, the future density in
physical coordinates is the same $M_t$ as before the restart: its
conformal radius and tip--target angle depend only on the same marked
physical domain.  There were no earlier $e_N$-approaches in
$Q$.  The normalized tail \eqref{mesh:normalized-tail}, applied in
$A_Q$ at scale $2l$, with the time cutoff deleting contributions,
therefore gives
\[
 \mathbb P\{\nu_N(Q)>u l^d\mid\mathcal F_{\sigma_Q}\}
 \le C\exp(-cu^{2/p})
\]
on $E_Q$, for $u\ge u_0$, uniformly in the final mesh $e_N\le l$;
here $u_0$ is enlarged by a fixed factor to account for the radius $2l$.
The case settled by the deterministic bound in the derivation of
\eqref{mesh:normalized-tail} includes meshes too coarse for its
intermediate test.  Provided $D(\log k_0)^{p/2}\ge u_0$, integrate
this tail above
$u=D(\log k)^{p/2}$, including the usual threshold-times-tail term.
Any polynomial prefactor is absorbed by reducing $c$.  Multiplication
by the entrance probability $Cl^p$ proves
\eqref{retain:bad-square}.

Let $E_N$ be the union of all excessive squares at levels
$k_0,\ldots,N$, and retain the measure
\[
 \widehat\nu_N=\mathbf1_{E_N^c}\nu_N.
\]
There are at most $Cl^{-2}$ squares meeting $K$ at level $k$.
Since $p+d=2$, summation of \eqref{retain:bad-square} gives
\[
 \mathbb E\bigl[\nu_N(K)-\widehat\nu_N(K)\bigr]
 \le C\sum_{k\ge k_0}k^{-\beta}.
\]
Choose $D$ so that $\beta>1$, then choose $k_0$ large enough to meet
the geometric and threshold requirements above and make this sum at most $m_0/2$.
It follows from \eqref{retain:first-second} that
\[
 \mathbb E\widehat\nu_N(K)\ge m_0/2,
 \qquad\mathbb E\widehat\nu_N(K)^2\le m_2.
\]
Paley--Zygmund therefore supplies deterministic $a_*,p_*>0$ such
that
\begin{equation}
 \mathbb P\{\widehat\nu_N(K)\ge a_*\}\ge p_*
 \quad\text{for all sufficiently large }N.
 \label{retain:positive-event}
\end{equation}

At every tested level each square has retained mass at most $H_k$:
an excessive square has retained mass zero, and otherwise restriction
can only reduce its mass.  In particular the finitely many squares
at level $k_0$ meeting $K$ give a deterministic bound on the total
retained mass, uniform in $N$.

\subsection{A supported weak limit and its ball bound}

The event in \eqref{retain:positive-event} occurs infinitely often
with probability at least $p_*$.  Indeed each union over $N\ge n$
has probability at least $p_*$, and these unions decrease to the
limsup event.  On that event choose a subsequence on which the mass
is at least $a_*$.  Measures on the fixed compact disk $K$ with the
deterministic total-mass bound are weakly sequentially compact.
Thus a further subsequence converges weakly to a measure $\nu$ with
$\nu(K)\ge a_*$.  No measurable selection of this subsequence is
needed: the pathwise existence of such a measure will imply a
measurable positive-Hausdorff-measure event.

The support of $\nu$ lies in $K\cap\gamma([0,T])$.  Every
$\widehat\nu_N$ is carried by the closed $e_N$-neighborhood of that
compact trace, and a compact set disjoint from the trace has an open
neighborhood disjoint from all these carriers for large $N$.

For completeness, weak convergence causes no difficulty at dyadic
boundaries.  Given a sufficiently small $r$, choose a fixed dyadic
side length $l$ with $r\le l<2r$.  For all sufficiently large $N$
this level was tested, and every open ball $B(x,r)$ meets only a
bounded number of its half-open squares.  Therefore
\[
 \widehat\nu_N(B(x,r))\le C D l^d(\log\log(1/l))^{p/2}
 \le C' h(r).
\]
Portmanteau for this open ball gives
\begin{equation}
 \nu(B(x,r))\le C'h(r).
 \label{retain:ball}
\end{equation}
This argument works for every $x,r$ on each path with a weak limit;
it does not assume zero mass on grid boundaries.  It also shows that
$\nu$ has no atoms.

Any set of positive diameter $r$ is contained in an open ball of
radius at most $2r$, centered at one of its points.  The bounded
scaling property of $h$ and \eqref{retain:ball} consequently imply
$\nu^*(U)\le C h(\operatorname{diam}U)$ for all sufficiently small
sets, where $\nu^*$ is the outer measure associated with $\nu$;
singleton sets have zero mass.  Countable subadditivity over an arbitrary
countable cover of the trace proves
\[
 \mathcal H^h(K\cap\gamma([0,T]))\ge \nu(K)/C>0.
\]
We have proved a fixed positive probability of positive Hausdorff
measure on a compact interior piece, in every prescribed nonzero
initial time interval after scaling.

\subsection{Restarting in every positive-time interval}\label{retain:positivity}

Fix $0\le s<t<\infty$ and divide $[s,t]$ into $n$ deterministic
capacity intervals.  At a left endpoint $t_i$, the hydrodynamic map
$g_{t_i}-\sqrt\kappa B_{t_i}$ turns the future into an ordinary
chordal half-plane evolution driven by
$\sqrt\kappa(B_{t_i+u}-B_{t_i})$.  Its capacity time is exactly $u$:
hydrodynamic composition adds the coefficient $2u$.  Thus no random
time change is hidden in this subdivision.

Apply the preceding scaled interior experiment in each fresh
half-plane interval.  Its success probability is at least $p_*$,
conditionally on the preceding past, with $p_*$ independent of the
interval length.  Success provides a compact subset strictly inside
the half-plane.  The inverse map is conformal in a neighborhood of
this compact subset.  A finite cover by interior disks gives upper
and lower metric distortion bounds locally, and hence preserves
positivity of $\mathcal H^h$ by the bounded scaling property of $h$.
One may restrict to one disk carrying positive measure if needed.
The mapped compact subset belongs to the original trace during that
same capacity interval.

Consequently failure of positivity on the whole $[s,t]$ implies
failure of all $n$ experiments, and has probability at most
$(1-p_*)^n$.  Letting $n\to\infty$ shows positivity almost surely for
this fixed interval.  The event is measurable because the trace is
a random compact set and Hausdorff measure for a continuous gauge
can be computed using the countable family of finite covers by
finite unions of rational open balls, with arbitrarily small
enlargements.  Intersecting over rational $0<s<t$ gives a single
probability-one event on which every real $0<s<t<\infty$ has
positive measure: each such interval contains a nontrivial rational
subinterval.  The event is for the one fixed value of $\kappa$.

\subsection{A prelimit lower-tail statement for ordinary sausages}

The first and second moments also give an input for the later ribbon
argument without passing to a limiting measure.  For a compact time
interval $J$ write
\[
 X_e(J)=e^{-p}\operatorname{Area}
  \{z\in\mathbb C:\operatorname{dist}(z,\gamma(J))<e\}.
\]

\begin{proposition}[Small sausage values are unlikely]
\label{retain:sausage}
For every deterministic $J=[s,t]$ with $0\le s<t<\infty$,
\begin{equation}
 \lim_{a\downarrow0}\limsup_{e\downarrow0}
 \mathbb P\{X_e(J)<a\}=0.
 \label{retain:sausage-tail}
\end{equation}
Equivalently, for every $\varepsilon>0$ there are $a,e_0>0$ such
that $\mathbb P\{X_e(J)<a\}<\varepsilon$ for all $0<e<e_0$.
The same conclusion can be required simultaneously for any fixed
finite list of such intervals, with total exceptional probability
as small as desired.
\end{proposition}

\begin{proof}
In the initial half-plane experiment, Paley--Zygmund applied directly
to \eqref{retain:first-second} gives constants $a_{\rm s},p_{\rm s}>0$
such that $\mathbb P\{\nu_{e/2}^T(K)>a_{\rm s}\}\ge p_{\rm s}$
for all sufficiently small $e$.  Since
$W_{e/2}(z)\le Ce^{-p}$ and the density is supported within distance
$e/2$ of the trace up to $T$, this implies a positive lower bound on
$e^{-p}$ times its sausage area inside $K$.  Use a slightly enlarged
compact disk $K^+\Subset\mathbb H$ to ensure that all witnessing
trace points lie in $K^+$ for small $e$.  Scaling gives this
experiment on any prescribed positive-length half-plane slot,
with the same lower bound $p_{\rm s}$ on its success probability.

Divide $J$ into $n$ deterministic slots and let $F_i$ be the inverse
hydrodynamic map at the start of slot $i$.  The corresponding test
disks $K_i\Subset K_i^+\Subset\mathbb H$ are deterministic after
scaling by that slot's length.  There are almost surely finite upper
and strictly positive lower bounds on $|F_i'|$ throughout $K_i^+$,
for each $i$.  Because the list is finite, for every $\eta>0$ we can
choose deterministic $0<b<L<\infty$ such that, except on an event
of probability at most $\eta$, all these derivatives lie between
$b$ and $L$.  Enlarge $L$ to include the local Lipschitz constants on
the relevant compact disks.  No independence between this event and
the slot experiments is asserted or needed.

On that event take each half-plane mesh to be $ce$, where the
deterministic $c>0$ is small enough that image distances less than
$ce$ become physical distances less than $e$.  If its slot
experiment succeeds, conformal area change, with Jacobian at least
$b^2$, gives $X_e(J)\ge a_{n,\eta}$, where $a_{n,\eta}>0$ is
deterministic.  This conclusion holds for all sufficiently small
$e$, with a deterministic cutoff depending on the finite list.
The half-plane experiments have conditional success probability at
least $p_{\rm s}$, so their simultaneous failure has probability at most
$(1-p_{\rm s})^n$.  Hence
\[
 \limsup_{e\downarrow0}\mathbb P\{X_e(J)<a_{n,\eta}\}
 \le\eta+(1-p_{\rm s})^n.
\]
First choose $n$ large and then $\eta$ small.  Monotonicity in the
threshold $a$ proves \eqref{retain:sausage-tail}.  For a fixed finite
list, apply this result to each interval and take the minimum
threshold and cutoff, followed by a union bound.
\end{proof}

Only deterministic intervals are asserted in
Proposition~\ref{retain:sausage}.  Its use for a random passage first
selects, by countable exhaustion, a finite list of rational closed
intervals and compact interior neighborhoods with positive margins.
The proposition then removes an arbitrarily small unconditional
exceptional event.  On a passage whose selected segment lies inside
one of those neighborhoods, sufficiently fine sausages of the segment
are counted there.  No conditional sausage inequality at an arbitrary
random interval is required.
